% last page to be written
% -> Most important aspects of the work: essential and significant
% -> Context; Relevance and Objectives; Methodology Used; Results and Conclusions
% -> Keywords

%% abstract.tex: abstract in PT and EN  (FEUP regulations)
%% -------------------------------------------------------
\chapter*{Resumo}
%

Com o uso crescente de modelos de linguagem de grande escala (LLMs) em contextos de tomada de decisão, torna-se essencial compreender os seus modos de falha e avaliar a sua performance em cenários complexos. A negociação é um contexto particularmente adequado para avaliar estes modelos, uma vez que se trata de um cenário exigente que requer interpretar as intenções de terceiros, planear a longo prazo e agir sob incerteza. A \textit{NegotiationArena} é uma framework de código aberto que disponibiliza um \textit{benchmark} desafiante para avaliar as capacidades de negociação dos LLMs em três cenários de referência: negociação de preços, troca de recursos e o Jogo do Ultimato. A avaliação da \textit{NegotiationArena} limitou-se a analisar modelos proprietários, alguns dos quais, entretanto, foram descontinuados. Nesta dissertação, avaliamos nove modelos de pesos abertos de três famílias (Qwen, Gemma e Mistral) e analisamos o impacto de duas técnicas aplicadas em tempo de inferência: o autoaperfeiçoamento, em que o agente revê e reformula a sua jogada antes de a submeter, e a negociação em equipa, em que três modelos deliberam para tomar uma decisão conjunta.

Os modelos de pesos abertos são capazes de negociar nos três cenários e reproduzem padrões anteriormente descritos para os modelos proprietários, como a ancoragem de preços e as vantagens posicionais. Identificamos alguma irracionalidade, incluindo rejeições que destroem valor para ambas as partes e propostas abaixo do custo de produção.

Embora haja uma grande variabilidade de resultados, o que era de esperar tendo em conta o número de variáveis, foi possível observar certos padrões. O Qwen é a família de modelos mais consistente, vencendo entre 63\% e 67\% dos jogos, seguido do Gemma e do Mistral. A introdução de duas \textit{personas}, a desesperada (\textit{desperate}) e a manipuladora (\textit{cunning}), que aumentaram o \textit{payoff} e a taxa de vitória do GPT-4 no artigo da \textit{NegotiationArena}, não têm o mesmo efeito nos modelos de pesos abertos; o seu impacto depende do jogo e do modelo. As duas técnicas aplicadas em tempo de inferência introduzidas nesta dissertação não produziram ganhos suficientes que justifiquem o seu custo. O autoaperfeiçoamento tem um efeito assimétrico, ajudando nas posições estruturalmente desfavorecidas e penalizando as que já são fortes. A negociação em equipa deixa a taxa de vitória praticamente inalterada e produz apenas ganhos de \textit{payoff} localizados. A nossa análise revela limitações dos LLMs de pesos abertos recentes em cenários de negociação desafiantes e aponta direções para melhorias em trabalhos futuros.

\vspace{\fill}
\begin{flushleft}
{\Large \textbf{Palavras-chave}:} 
Modelos de Linguagem de Grande Escala \(\bullet\) 
Negociação \(\bullet\)
Técnicas em Tempo de Inferência \(\bullet\)
Colaboração Multiagente

\end{flushleft}
\vspace*{24mm}

\acresetall %% to reset the acronym usage

\chapter*{Abstract}
% Here goes the abstract written in English. It should say the same.

With the increasing use of large-scale language models (LLMs) in decision-making contexts, it becomes essential to understand their failure modes and evaluate their performance in complex scenarios. Negotiation is a particularly suitable context for evaluating these models, as it is a demanding scenario that requires interpreting the intentions of third parties, long-term planning, and acting under uncertainty. \textit{NegotiationArena} is an open-source framework that provides a challenging benchmark to assess the negotiation capabilities of LLMs in three reference scenarios: price negotiation, resource exchange, and the Ultimatum Game. The \textit{NegotiationArena} evaluation was limited to analyzing proprietary models, some of which have since been discontinued. In this dissertation, we evaluate nine open-weight models from three families (Qwen, Gemma, and Mistral) and analyse the impact of two inference-time techniques: \textit{Self-Refine}, where the agent reviews and reformulates its move before submitting it, and team negotiation, where three models deliberate to make a joint decision.

Open-weight models can negotiate in all three scenarios and reproduce some patterns previously described for proprietary models, such as price anchoring and positional advantages. We identify some irrationality, including rejections that destroy value for both parties and bids below the cost of production.

Although there is considerable variability in results, as expected given the number of variables, certain patterns emerged. Qwen is the most consistent family of models, winning between 63\% and 67\% of the games, followed by Gemma and Mistral. The \textit{desperate} and \textit{cunning} personas, which increased GPT-4's \textit{payoff} and \textit{win rate} in the \textit{NegotiationArena} article, do not have the same effect on open-weight models; their impact depends on the game and the model. The two inference-time techniques introduced in this dissertation did not produce sufficient gains to justify their cost. \textit{Self-Refine} has an asymmetrical effect, helping structurally disadvantaged positions and penalising those that are already strong. Team negotiation leaves the win rate virtually unchanged and produces only localised \textit{payoff} gains. Our analysis reveals limitations of recent open-weight LLMs in challenging negotiation scenarios and points to directions for future work improvements.

\vspace{\fill}
\begin{flushleft}
{\Large \textbf{Keywords}:} 
Large Language Models \(\bullet\) 
Negotiation \(\bullet\)
Inference-Time Techniques \(\bullet\)
Multi-Agent Collaboration
\end{flushleft}
\vspace*{24mm}

\acresetall %% to reset the acronym usage
